% Chapter 3
\chapter{Methodology and Requirements Planning}
\label{chap:planning}

This chapter presents how the dissertation project will be carried out during execution (January--June 2026). Building on the problem framing in Chapter~\ref{chap:introduction} and state of the art in Chapter~\ref{chap:literature}, the focus is on methodological approach, planned artifacts, evaluation activities, skills development, schedule, and risk management.

\section{Objectives and Scope Alignment}

The project objectives and scope defined in Chapter~\ref{chap:introduction} remain aligned with findings synthesized in Chapter~\ref{chap:literature}. No scope changes are planned. The work focuses on execution strategy to deliver planned artifacts and produce defensible evidence.

\section{Requirements Definition}
\label{sec:ch3_requirements}

This section consolidates the functional and non-functional requirements that guide architecture and implementation decisions.

\textbf{Functional requirements} include: collection of metrics, logs, and traces across core \ac{CICD} services; centralized visualization; and support for cross-service incident analysis.

\textbf{Non-functional requirements} include: operation in heterogeneous environments (Linux and Windows), reproducible deployment procedures, acceptable instrumentation overhead, and operational maintainability for infrastructure teams.

These requirements are used as acceptance criteria during design reviews and as reference points for evaluation planning in subsequent chapters.

\section{Methodological Approach}
\label{sec:ch3_method}

The project is planned as an engineering-oriented research effort centered on the design, implementation, and evaluation of an observability solution for heterogeneous containerized environments. The chosen approach is pragmatic and iterative: requirements and platform characteristics inform successive design decisions, which are validated through controlled experimentation and practitioner-oriented feedback.

This section describes the methodological structure guiding the dissertation execution. The subsections specify the concrete artifacts to be produced, the development process ensuring traceability from requirements to design decisions, and the evaluation strategy for generating defensible empirical evidence. Together, these components establish how the research moves from problem framing and literature synthesis to implemented solutions and validated outcomes.

\subsection{Planned Artifacts}
\label{sec:ch3_artifacts}

The project will produce a set of artifacts intended to be directly auditable and reusable:
\begin{itemize}
    \item \textbf{A1 -- Reference architecture:} a documented architecture describing components, deployment topology, data flows, and correlation mechanisms for heterogeneous containerized environments, including comparative analysis of alternative instrumentation patterns and documented trade-off evaluation.
    \item \textbf{A2 -- Prototype implementation:} a working prototype implementing the reference architecture in a staging environment, including configuration and operational procedures.
    \item \textbf{A3 -- Automation and reproducibility assets:} scripts and configuration files to provision, deploy, and execute the evaluation campaign in a repeatable manner.
    \item \textbf{A4 -- Evaluation protocol and dataset:} an experimental plan (scenarios, workloads, baselines, metrics) and the resulting collected data, with clear provenance.
    \item \textbf{A5 -- Documentation package:} technical documentation for the architecture, deployment, and evaluation procedures, prepared to support both thesis writing and oral defense.
\end{itemize}

\subsection{Development Process and Traceability}
\label{sec:ch3_process}

The development process ensures traceability from requirements to design decisions and from design decisions to evaluation evidence:
\begin{itemize}
    \item \textbf{Requirements-to-design mapping:} platform characteristics and operational requirements (e.g., heterogeneous execution environments, cross-platform support) are mapped to explicit architectural and implementation decisions.
    \item \textbf{Iterative refinement:} the architecture and prototype are developed in short iterations, each concluding with a review of the produced artifact increment and the update of open issues.
    \item \textbf{Configuration management:} all artifacts are maintained under version control, including deployment manifests, configuration, and evaluation scripts.
    \item \textbf{Acceptance criteria:} each major deliverable is associated with objective acceptance criteria (e.g., successful deployment steps, reproducible execution of an experiment, completeness of documentation).
\end{itemize}

\subsection{Planned Evaluation Strategy}
\label{sec:ch3_evaluation}

Evaluation is planned to produce evidence that the proposed solution provides measurable utility when compared with a defined baseline, and to characterize trade-offs between alternative instrumentation approaches. The evaluation design is bounded to what can be executed within the dissertation schedule and the available infrastructure.

		\textbf{Evaluation dimensions.} The evaluation covers:
\begin{itemize}
    \item \textbf{Architectural alternatives:} comparative analysis of different instrumentation patterns (sidecar, privileged containers, host-level agents) against defined criteria including implementation complexity, resource overhead, and diagnostic capability.
    \item \textbf{Operational overhead:} resource consumption and performance overhead introduced by the solution (measured using a defined workload and baseline).
    \item \textbf{Observability usefulness:} the extent to which the collected telemetry supports diagnosing representative failure and degradation scenarios.
    \item \textbf{Maintainability and operability:} the clarity of configuration, deployment, and troubleshooting procedures as captured in the documentation.
\end{itemize}

		\textbf{Experimental design (planning level).} The experimental campaign is structured as follows:
\begin{itemize}
    \item \textbf{Baselines:} at least one baseline configuration is defined to enable comparative measurements (e.g., no additional observability components, or minimal instrumentation).
    \item \textbf{Workloads and scenarios:} a small set of representative workloads and fault/degradation scenarios are specified to exercise logging, metrics, and tracing paths.
    \item \textbf{Data collection:} procedures are documented to collect resource metrics, request/trace metadata, and logs, ensuring consistent time windows and environment configuration.
    \item \textbf{Analysis:} the analysis plan defines how measurements are summarized and compared, and how qualitative observations are recorded.
\end{itemize}

		\textbf{Practitioner feedback (planning level).} Where access and availability permit, the project includes structured feedback from practitioners (e.g., hosting organization stakeholders) using an interview or review protocol. The protocol defines eligibility, consent, the question guide, and how responses are summarized (e.g., thematic coding at a coarse granularity) to support defensible qualitative claims.

\section{Skills Development Plan (PREPD)}
\label{sec:ch3_skills}

This section defines a skills development plan aligned with the project phases. The plan identifies the competencies required for the dissertation, the current diagnosis, development actions, and objective assessment criteria.

The competency framework organizes required capabilities into three domains: observability engineering in constrained environments, experimental evaluation with attention to validity, and technical communication for both written and oral presentation. The subsection that follows diagnoses current and target competency levels, specifies planned development actions mapped to project phases, and defines objective criteria for assessing whether competency targets have been achieved.

\subsection{Competency Diagnosis}

\begin{table}[htbp]
\centering
\caption{Competency Diagnosis and Targets}
\label{tab:competencies}
\footnotesize
\begin{tabular}{@{}p{0.22\textwidth}p{0.25\textwidth}p{0.25\textwidth}p{0.22\textwidth}@{}}
		\textbf{Competency} & \textbf{Current Level} & \textbf{Target Level} & \textbf{Relevance} \\
\midrule
		\textbf{C1: Observability engineering in containers} &
Intermediate: operational familiarity with containers and basic telemetry tooling &
Advanced: design and deploy a production-oriented telemetry stack; evaluate alternative instrumentation patterns &
Core competency to implement A1--A3 \\
\addlinespace
		\textbf{C2: Experimental evaluation and validity} &
Foundational: understanding of evaluation concepts; limited campaign execution &
Intermediate: define baselines, run experiments, analyze results, discuss threats &
Required to produce A4 and defensible evidence \\
\addlinespace
		\textbf{C3: Technical communication and stakeholder interaction} &
Intermediate: academic writing; limited structured practitioner engagement &
Advanced: concise documentation, structured feedback collection, synthesis &
Required for A5 and oral defense readiness \\
\bottomrule
\end{tabular}
\end{table}


\begin{table}[htbp]
\centering
\caption{Skills Development Actions Mapped to Project Phases}
\label{tab:development}
\footnotesize
\begin{tabular}{@{}p{0.14\textwidth}p{0.48\textwidth}p{0.30\textwidth}@{}}
		\toprule
		\textbf{Comp.} & \textbf{Planned Development Actions (by phase)} & \textbf{Assessment Criteria} \\
\midrule
		\textbf{C1} &
		\textbf{P1--P2:} deep-dive on candidate tooling and deployment patterns; comparative evaluation of instrumentation approaches; implement staging deployment; document operational procedures.\newline
		\textbf{P3:} iterate architecture and prototype based on deployment feedback and evaluation of alternatives. &
Prototype deploys from documented procedure; architecture decisions include trade-off analysis; operational runbook is complete; alternative patterns are evaluated and documented. \\
\addlinespace
		\textbf{C2} &
		\textbf{P2:} define evaluation protocol (metrics, baselines, scenarios, data collection).\newline
		\textbf{P3--P4:} execute measurement campaign; perform comparative analysis; document threats to validity. &
Evaluation protocol is reproducible; dataset provenance is documented; analysis includes baselines and validity discussion. \\
\addlinespace
		\textbf{C3} &
		\textbf{P1:} consolidate PREPD writing and structure.\newline
		\textbf{P3--P5:} produce technical documentation; prepare interview/review protocol; synthesize practitioner feedback where applicable. &
Documentation supports independent replication; feedback protocol is structured; writing and oral preparation are consistent and coherent. \\
\bottomrule
\end{tabular}
\end{table}

\section{Project Planning and Schedule}
\label{sec:ch3_planning}

The project plan is organized into phases with explicit deliverables and dependencies. The sequencing is designed to reduce risk by validating feasibility early (technology selection and staging deployment) before committing to extended evaluation work.

This section presents the temporal organization of dissertation activities from January through June 2026. The subsections describe the phase structure with associated activities and deliverables, define milestones and their dependencies to support periodic progress assessment, and provide a visual timeline through a Gantt chart. This structure enables systematic tracking of progress and early identification of schedule risks.

\subsection{Phases, Activities, and Deliverables}

\begin{table}[htbp]
\centering
\caption{Planned Phases, Activities, and Deliverables (January--June 2026)}
\label{tab:phases}
\footnotesize
\begin{tabular}{@{}p{0.10\textwidth}p{0.20\textwidth}p{0.42\textwidth}p{0.22\textwidth}@{}}
		\toprule
		\textbf{Phase} & \textbf{Period} & \textbf{Key Activities} & \textbf{Main Deliverables} \\
\midrule
		\textbf{P1} & Jan &
Project setup; consolidation of requirements and constraints; definition of acceptance criteria and logs (issues/risks). &
Execution plan; project logs; refined acceptance criteria. \\
\addlinespace
		\textbf{P2} & Jan--Feb &
Technology assessment and selection; proof-of-concept validations; definition of the architecture decision record structure. &
Selection matrix and justification; initial architecture decisions. \\
\addlinespace
		\textbf{P3} & Feb--Mar &
Reference architecture specification; prototype implementation; staging deployment; automation and reproducibility assets. &
A1--A3 (architecture, prototype, automation); deployment runbook. \\
\addlinespace
		\textbf{P4} & Apr--May &
Evaluation campaign execution: baseline definition, scenarios, measurements, data collection, preliminary analysis. &
A4 (evaluation protocol and dataset); preliminary analysis notes. \\
\addlinespace
		\textbf{P5} & May--Jun &
Consolidated analysis; practitioner feedback (when feasible); thesis writing and revision cycles; defense preparation. &
A5 (documentation package); completed thesis manuscript. \\
\bottomrule
\end{tabular}
\end{table}

\subsection{Milestones and Dependencies}

Milestones are defined to support periodic assessment and early detection of schedule risk. Key dependencies are: (i) technology selection before architecture freeze, (ii) staging deployment before extended evaluation, and (iii) availability of infrastructure and stakeholders before practitioner feedback.

\begin{table}[htbp]
\centering
\caption{Planned Milestones and Acceptance Focus}
\label{tab:milestones}
\footnotesize
\begin{tabular}{@{}p{0.10\textwidth}p{0.22\textwidth}p{0.26\textwidth}p{0.34\textwidth}@{}}
		\toprule
		\textbf{ID} & \textbf{Planned Date} & \textbf{Milestone} & \textbf{Acceptance Focus} \\
\midrule
		\textbf{M1} & End Jan & Technology selection completed & Selection rationale documented; fallbacks identified; feasibility validated via proof-of-concept. \\
\addlinespace
		\textbf{M2} & Mid Mar & Staging prototype deployed & Deployment is repeatable; core telemetry pipeline works end-to-end; runbook is usable. \\
\addlinespace
		\textbf{M3} & End Apr & Evaluation protocol frozen & Baselines and scenarios defined; collection procedures documented; validity considerations recorded. \\
\addlinespace
		\textbf{M4} & End May & Measurement campaign executed & Dataset collected with provenance; key analyses executed; issues and limitations captured. \\
\addlinespace
		\textbf{M5} & Mid Jun & Thesis ready for submission & Chapters integrated; evidence traceable to artifacts; revisions incorporated; defense preparation completed. \\
\bottomrule
\end{tabular}
\end{table}

\subsection{Gantt Chart}

\begin{figure}[htbp]
\centering
\footnotesize
\begin{ganttchart}[
	hgrid,
	vgrid,
	x unit=0.32cm,
	y unit chart=0.55cm,
	bar height=0.55,
	group height=0.7,
	milestone height=0.6,
	milestone label font=\footnotesize,
	bar label font=\footnotesize,
	group label font=\footnotesize,
]{1}{26}

\gantttitle{2026}{26}\\
\gantttitle{Jan}{4}\gantttitle{Feb}{4}\gantttitle{Mar}{5}\gantttitle{Apr}{4}\gantttitle{May}{5}\gantttitle{Jun}{4}\\

\ganttgroup{Project execution (Jan--Jun 2026)}{1}{26}\\

\ganttbar[name=p1]{P1: Setup and planning}{1}{4}\\
\ganttbar[name=p2]{P2: Tooling assessment and selection}{3}{8}\\
\ganttbar[name=p3]{P3: Architecture and prototype}{5}{13}\\
\ganttbar[name=p4]{P4: Evaluation campaign}{14}{22}\\
\ganttbar[name=p5]{P5: Analysis, writing, defense prep}{19}{26}\\

\ganttmilestone[name=m1]{M1: Tech selection complete}{4}\\
\ganttmilestone[name=m2]{M2: Staging prototype deployed}{11}\\
\ganttmilestone[name=m3]{M3: Eval protocol frozen}{17}\\
\ganttmilestone[name=m4]{M4: Measurements executed}{22}\\
\ganttmilestone[name=m5]{M5: Thesis submission-ready}{25}

\ganttlink{p2}{m1}
\ganttlink{p3}{m2}
\ganttlink{p4}{m3}
\ganttlink{p4}{m4}
\ganttlink{p5}{m5}

\end{ganttchart}
\caption{Project schedule overview (Gantt chart; January--June 2026)}
\label{fig:gantt}
\end{figure}

\section{Project Monitoring and Management}
\label{sec:ch3_monitoring}

Project monitoring is structured to ensure timely control of scope, schedule, and quality:
\begin{itemize}
    \item \textbf{Governance and roles:} responsibilities are split between the student (execution and reporting), the academic advisor (scientific guidance and quality control), and the hosting institution supervisor (infrastructure coordination and practitioner feedback facilitation when applicable).
    \item \textbf{Cadence:} meetings are held regularly (e.g., bi-weekly with the academic advisor; weekly or as-needed with the hosting supervisor) to review progress against the plan and to unblock dependencies.
    \item \textbf{Progress tracking:} a lightweight project log is maintained, including tasks, deliverables, open issues, decisions, and risk status.
    \item \textbf{Change control:} any scope adjustment is documented with rationale, impact analysis (schedule/risk), and approval by the advisor.
    \item \textbf{Quality control:} each milestone includes an acceptance review focused on reproducibility, completeness of documentation, and alignment with constraints.
\end{itemize}

\section{Risk Management}
\label{sec:ch3_risks}

Risk management is continuous and integrated with monitoring. Risks are identified early, assessed by probability and impact, tracked through indicators, and mitigated through preventive actions. When mitigation is insufficient, a contingency plan is activated to preserve the ability to complete the dissertation within the time constraints.

To keep the assessment consistent and auditable, this project uses a lightweight qualitative scale for both \,\textit{Impact} and \,\textit{Probability}. Impact captures the consequence to delivery (scope, schedule, or validity of results); probability captures the likelihood of occurrence within the January--June 2026 execution window. Monitoring indicators are defined as observable signals with thresholds (where feasible) that trigger mitigation or contingency actions.

\begin{table}[htbp]
\centering
\caption{Risk Rating Scales and Measurement Approach}
\label{tab:risk_scales}
\footnotesize
\begin{tabular}{@{}p{0.14\textwidth}p{0.18\textwidth}p{0.60\textwidth}@{}}
	\toprule
	\textbf{Dimension} & \textbf{Class} & \textbf{Operational definition (how it is measured)} \\
\midrule
	\textbf{Impact} & H & Prevents delivering a core artifact (A1--A4) or forces a major scope reduction; or causes a schedule slip \textgreater{} 2 weeks on a milestone. \\
 & M & Degrades evidence quality or requires rework; or causes a schedule slip of \textasciitilde{} 1--2 weeks. \\
 & L & Localized issue with limited rework; schedule impact \textless{} 1 week. \\
\addlinespace
	\textbf{Probability} & H & Likely (\textgreater{} 60\%): already observed signals or strong prior evidence in similar environments. \\
 & M & Possible (20--60\%): plausible given constraints; may occur depending on approvals/changes. \\
 & L & Unlikely (\textless{} 20\%): would require unusual conditions or multiple contributing events. \\
\bottomrule
\end{tabular}
\end{table}

\begin{table}[htbp]
\centering
\caption{Risk Register (technical, schedule, and external risks)}
\label{tab:risks}
\scriptsize
\begin{tabular}{@{}p{0.06\textwidth}p{0.11\textwidth}p{0.26\textwidth}p{0.07\textwidth}p{0.07\textwidth}p{0.18\textwidth}p{0.18\textwidth}@{}}
	\toprule
		\textbf{ID} & \textbf{Type} & \textbf{Risk Description} & \textbf{Impact} & \textbf{Prob.} & \textbf{Monitoring Indicators} & \textbf{Mitigation / Contingency} \\
\midrule
R1 & Tech & Selected tooling or deployment pattern does not meet operational requirements, or is not reliable in the heterogeneous target environment. & H & M & Proof-of-concept fails to deploy \textgreater{}2 times; required capability missing; unstable telemetry pipeline. & Run early proof-of-concept; evaluate multiple architectural alternatives; keep fallback options; freeze a minimal viable architecture and reduce scope if needed. \\
\addlinespace
R2 & Tech & Telemetry pipeline quality and/or overhead is insufficient for the planned diagnosis scenarios. & H & M & Trace/log correlation breaks; missing identifiers; CPU/memory overhead exceeds agreed bounds; latency regression beyond baseline tolerance. & Validate correlation end-to-end early; tune configuration and sampling; prioritize fewer scenarios with stronger instrumentation; bound evaluation and document limitations. \\
\addlinespace
R3 & Schedule & Infrastructure access, deployment approvals, or environment instability delays evaluation activities. & H & M & Access not granted by planned date; blocked deployments; long coordination cycles; repeated environment outages. & Use staging as primary target; parallelize protocol design and writing; maintain schedule buffers; adjust campaign size to preserve M3--M5. \\
\addlinespace
R4 & Schedule & Writing and review iteration compresses due to technical delays, reducing dissertation quality. & H & M & Missed weekly writing checkpoint; advisor review windows missed; late integration of chapters/evidence. & Enforce weekly writing output; draft chapters incrementally; lock review dates; de-scope non-essential work to protect submission readiness. \\
\addlinespace
R5 & External & Limited stakeholder availability or ecosystem/policy changes reduce feedback feasibility or impose new constraints. & M & M & Low response rate; repeated rescheduling; new restrictions; deprecations impacting selected components. & Recruit early; allow asynchronous reviews; complement with advisor/expert review; choose conservative components; document assumptions and adapt plan. \\
\bottomrule
\end{tabular}
\end{table}

The risk register is reviewed at least monthly and at each milestone review. For high-impact risks, mitigation status and contingency readiness are reported in progress updates.
